\section*{TD 2 --- Préparation des données}
\addcontentsline{toc}{section}{TD 2 --- Préparation des données}

\subsection*{Exercice 1 --- Que faire des manquants ? \hfill \normalfont\footnotesize (8 min)}
Une table de 1\,000 étudiants présente 41 \% de valeurs manquantes sur le
revenu des parents, et 3 \% sur la moyenne générale.
\begin{enumerate}
  \item Proposer une décision pour chaque colonne, en justifiant.
  \item Pourquoi supprimer les lignes à 41 \% de manquants est-il risqué ?
  \item Quelle information ajouter pour améliorer l'imputation ?
\end{enumerate}

\subsection*{Exercice 2 --- Encodages \hfill \normalfont\footnotesize (8 min)}
Trois colonnes : \texttt{filiere} (informatique, télécom, réseaux),
\texttt{niveau} (faible, moyen, élevé), \texttt{code\_postal} (400 modalités).
Indiquer l'encodage retenu pour chacune et la raison.

\subsection*{Exercice 3 --- Effet de l'échelle \hfill \normalfont\footnotesize (10 min)}
Deux étudiants sont décrits par (ancienneté en années, revenu en FCFA) :
$A=(0,0)$ et $B=(3,\,4\,000)$.
\begin{enumerate}
  \item Calculer la distance euclidienne brute.
  \item Les écarts-types valent 1,5 an et 2\,000 FCFA. Recalculer la distance
        après standardisation.
  \item Que conclure sur l'usage de kNN sans normalisation ?
\end{enumerate}

\subsection*{Exercice 4 --- Détecter une valeur aberrante \hfill \normalfont\footnotesize (8 min)}
Série : $10,\,12,\,14,\,16,\,18,\,20,\,22,\,60$. On utilise la méthode de Tukey
(médiane de chaque moitié).
\begin{enumerate}
  \item Calculer $Q_1$, $Q_3$, l'IQR et les deux bornes.
  \item La valeur 60 est-elle aberrante ? Et la valeur 10 ?
  \item Faut-il la supprimer ? Que faut-il vérifier avant ?
\end{enumerate}

\subsection*{Exercice 5 --- Trouver la fuite dans le code \hfill \normalfont\footnotesize (10 min)}
\begin{lstlisting}
X_scaled = StandardScaler().fit_transform(X)
X_train, X_test, y_train, y_test = train_test_split(X_scaled, y, test_size=0.2)
imputer = SimpleImputer(strategy="mean").fit(X)
X_train = imputer.transform(X_train)
\end{lstlisting}
Deux lignes posent problème. Les identifier et réécrire le pipeline correctement.


\subsection*{Exercice 6 --- Encodage cyclique \hfill \normalfont\footnotesize (8 min)}
L'heure est encodée par $h_{\sin}=\sin(2\pi h/24)$ et $h_{\cos}=\cos(2\pi h/24)$.
\begin{enumerate}
  \item Calculer le couple $(\sin,\cos)$ pour 0 h, 6 h, 12 h et 18 h.
  \item Quelle est la distance euclidienne entre 23 h et 0 h ? Entre 0 h et 12 h ?
  \item Pourquoi un encodage entier de l'heure est-il trompeur ?
\end{enumerate}

\subsection*{Exercice 7 --- Imputation et dispersion \hfill \normalfont\footnotesize (7 min)}
On impute 25 \% de valeurs manquantes d'un revenu très asymétrique.
\begin{enumerate}
  \item L'imputation par la moyenne conserve-t-elle la moyenne ? La variance ?
  \item Que fait l'imputation par la médiane à la moyenne globale ?
  \item Quelle information perd-on dans les deux cas ?
\end{enumerate}

\subsection*{Exercice 8 --- Discrétiser ou non ? \hfill \normalfont\footnotesize (6 min)}
On découpe l'âge en quatre tranches.
\begin{enumerate}
  \item Citer deux avantages de ce choix.
  \item Citer deux inconvénients.
  \item Pour quel type de modèle la discrétisation est-elle la moins utile ?
\end{enumerate}

\subsection*{Exercice 9 --- La jointure qui double les lignes \hfill \normalfont\footnotesize (8 min)}
Une jointure fait passer de 400 étudiants à 752 lignes.
\begin{enumerate}
  \item Expliquer la cause probable.
  \item Quel effet sur le taux de réussite moyen ?
  \item Proposer deux corrections possibles.
\end{enumerate}

\newpage
